\chapter{Proprietà FLTL e verifica}

Le proprietà di vivacità sono condizionate all'assenza di un allarme
attivo, poiché l'evacuazione ha priorità sul servizio ordinario.
\begin{description}
  \item[\texttt{MUTUAL\_EXCLUSION}] dopo un ingresso, prima di un
    nuovo ingresso devono avvenire \texttt{exit} oppure
    \texttt{evacuate}:
    \[
      \Box\bigl(\mathit{enter}\Rightarrow
      \neg\mathit{enter}\ \mathsf{U}\
      (\mathit{exit}\vee\mathit{evacuate})\bigr).
    \]
  \item[\texttt{EVENTUAL\_ARRIVAL}] ogni chiamata è seguita dal
    servizio al piano dell'utente, se l'allarme viene risolto:
    \[
      \Box\bigl(\mathit{call}\Rightarrow
      \mathsf{F}\mathit{enter}\bigr).
    \]
  \item[\texttt{EVENTUAL\_DELIVERY}] ogni selezione è seguita
    dall'uscita oppure dall'evacuazione:
    \[
      \Box\bigl(\mathit{select}\Rightarrow
      \mathsf{F}(\mathit{exit}\vee\mathit{evacuate})\bigr).
    \]
  \item[\texttt{STARVATION\_FREEDOM}] un utente che ha effettuato una
    chiamata raggiunge prima o poi \texttt{enter}, assumendo fairness
    e risoluzione dell'allarme:
    \[
      \Box\bigl(\mathit{call}\Rightarrow
      \mathsf{F}\mathit{enter}\bigr).
    \]
\end{description}

In LTSA la proprietà di mutua esclusione può essere espressa come
processo parametrico:
\begin{verbatim}
property MUTUAL_EXCLUSION =
    (enter[h:FLOORS][u:USER_IDX] -> IN_CABIN),
IN_CABIN =
    (exit[h:FLOORS][u:USER_IDX][t:FLOORS] ->
        MUTUAL_EXCLUSION
    |evacuate[h:FLOORS][u:USER_IDX] ->
        MUTUAL_EXCLUSION).
\end{verbatim}
La stessa tecnica viene applicata alle proprietà di arrivo, consegna
e assenza di starvation, mantenendo gli indici per distinguere
utenti e piani.

\section{Esito dell'analisi}
La struttura del modello garantisce la mutua esclusione: \texttt{LIFT}
abilita \texttt{enter} soltanto nello stato \texttt{EMPTY} e porta la
cabina nello stato \texttt{FULL}. Le guardie del processo di coda
garantiscono inoltre che solo il ticket con \texttt{k == 0} possa
richiedere il servizio.

Le proprietà di vivacità richiedono l'ipotesi di fairness delle
transizioni abilitate e la risoluzione dell'allarme. Senza tale
ipotesi, una sequenza infinita di allarmi è ammessa dal modello e
rende impossibile garantire una consegna ordinaria. Il repository
contiene il modello FSP aggiornato; l'output numerico della verifica
LTSA non è tra gli artefatti presenti e va prodotto eseguendo le
proprietà sul composto \texttt{SYS}.